\chapter{执行语义与验收}
\label{ch:execution}

Certificate 恰含 forward 与 reverse。每个 Block 先执行有序 Steps，再执行一个 End。
两个根均从各自完整模型、空局部事实和空路径假设开始，共享同一个工作预算。
第一方向成功不会给第二方向增加事实，任一方向失败均不接受整个证书。
Block 证明完整输出构造；Derivation 只为 Cases 证明一条共同事实。

\section{引用解析与局部定义}
\label{sec:fact-numbering}

九种 Ref 的完整语义如下。编号均是相应逻辑空间的索引，而非文件偏移。
\begin{longtable}{p{0.22\textwidth}p{0.70\textwidth}}
\toprule
引用 & 解析结果与有效条件 \\
\midrule
\endhead
\code{VarLower}(i) & 当前 $I$ 的 $\mathrm{LE}(l_i-x_i)$；须存在有限下界。\\
\code{VarUpper}(i) & 当前 $I$ 的 $\mathrm{LE}(x_i-u_i)$；须存在有限上界。\\
\code{VarInt}(i) & 当前 $I$ 的 $\mathrm{INT}(x_i)$；须声明整数。\\
\code{RowLower}(j) & $\mathrm{LE}(L_j-a_j)$；须有有限下侧。\\
\code{RowUpper}(j) & $\mathrm{LE}(a_j-U_j)$；须有有限上侧。\\
\code{RowEq}(j) & $\mathrm{EQ}(a_j-L_j)$；须有相等的两有限侧。\\
\code{Local}(id) & $\Delta$ 中当前可见的已验定义；已 Drop、未定义或仅在兄弟可见时拒绝。\\
\code{Assumption}(k) & $\Gamma_k$；0 为最早加入的假设，越界拒绝。\\
\code{Exact} & 只对 checker 已确定的目标 proposition 做窄规则检查，不解析为普通事实。\\
\bottomrule
\end{longtable}

Lin、Round、EqPair、Cases、Plotkin 用显式 u32 \code{id} 发布一个局部事实。
ID 不须从零连续编号，也不须递增。
发布时该 ID 不得仍然活跃，且活跃数不得超过预算；所有前提、类型和表达式先通过检查。
Drop 按其数组顺序删除当前可见 ID，缺失 ID 或同一数组中重复删除已删 ID 均拒绝。
一个 ID 仅在删除后才能复用，新的定义有新的 proposition；引用只读当前定义。

\subsection{父子作用域}
\label{sec:branch-transactions}

Split 和 Cases 的两个子作用域继承父局部事实与既有假设。
子体可删除继承 ID，在该子体内复用 ID，并创建其他事实；这些变化不能改变父或兄弟。
两子体各附加本分支假设，退出时撤去该假设。
Cases 两子体都成功后，仅在父作用域发布其指定的一条共同结论；
子体的其他事实、删除或重定义均不向父发布。
Split 是端点，不存在需要合并的父续体。

lift 把当前活跃事实和路径假设带入其唯一子体。
旧坐标和旧行的完整前缀不变，因而旧事实仍有效；所有新行只在子体可见。
在分支内扩展不会给兄弟增加坐标或行。
实现可共享不可变父状态并保存子差量；这种存储选择不能改变上述可见性。

\section{目标命题检查与 Exact}
\label{sec:exact-check}

设 checker 要证明 $(K,e)$，通常引用 $r$ 解析为 $(K_r,e_r)$。
接受当且仅当 $e_r=e$ 精确成立，且 $K_r=K$ 或 $K_r=\code{Eq}$。
比例相等、仅在某个点取值相等、ID 相等及显示文本相等均不足。

Exact 不从证书读入命题；checker 先构造 $e$，再检查：
\begin{enumerate}
 \item EQ：$e$ 无变量项且常数为零。
 \item LE：$e$ 无变量项且常数不正。
 \item INT：常数为整数，每项系数为整数，且每个出现坐标声明整数。
\end{enumerate}
常数整数没有坐标前提。
非恒真的仿射不等式不能通过 Exact；无模型整数性支持的非恒定表达式也不能通过 Exact。
生成的菜单或 signed 二元坐标可满足 INT 的窄规则，
生成的 flow、interval 和 subflow 连续坐标不能。

Exact 可出现在 Return 检查、lift 前提和 DerivationEnd Witness。
它不能作为 Lin（即使乘子为零）、Round、EqPair、Plotkin、整数 Split/Cases
或任何 Absurd 的原始事实操作数。
一个已有 Eq 可做目标检查的种类弱化，但不赋予原始操作数自动种类转换。

\section{步骤序列与穷尽分类}

步骤顺序执行；后一步仅能引用当前有效的模型事实、假设与前面活跃定义。
各条推理规则见第\ref{ch:instructions}章。
六种 Step 的命名空间与八种 End、两种 DerivationEnd 的命名空间独立。
同一个数字 tag 在不同位置不表示同一操作。

Split 与 Cases 都先验证规范表达式 $e$。
未提交整数引用时，两个分支假设为
\[
 \mathrm{LE}(e),\qquad\mathrm{LE}(-e).
\]
它们覆盖全实数域，$e=0$ 的边界同时属于两支。
提交整数引用时，须解析出\emph{种类确为 Int}且表达式恰为 $e$ 的事实，
然后两假设为
\[
 \mathrm{LE}(e),\qquad\mathrm{LE}(1-e).
\]
整数域中分别表示 $e\le0$ 与 $e\ge1$，穷尽且互斥。
不能用 Exact 或 Eq 原始引用冒充该 Int 前提。

\subsection{Cases 的共同事实}

Cases 在父作用域指定 ID、Kind 和 expression；两 Derivation
分别在穷尽分支上证明这同一 proposition，或用普通 Absurd 证明该支为空。
Witness 使用目标命题检查规则，因此允许 Exact 以及同表达式的 Eq 种类弱化。
两支不能各自选择不同 expression，再要求 checker 猜测共同结论。
Derivation 没有 map、Return、lift。
内部 Steps 仍可递归 Cases；成功后父得到唯一指定事实。

\subsection{Split 的完整子证明}

每个 Split 子 Block 都必须以该上下文允许的端点成功完成。
两支 Return 可以给出不同完整输出向量；Absurd 可关闭已证明不可能的支。
不能把一个分支只验了变量界、另一分支只验了行的证明合成一个 Return。
判为不可能的支仍须先证明合法矛盾；kernel 不按猜测的可达性跳过证明。


\section{完整 Return 与义务生成}
\label{sec:mapping-state}
\label{sec:obligation-generation}
map 长度恰等于输出变量数，表达式使用当前输入坐标并须规范。
对同一个完整向量 $Y$，checks 依次证明：
\begin{enumerate}
\item 输出变量按坐标序：先有限下界 $l_j-Y_j\le0$，再有限上界 $Y_j-u_j\le0$，
然后对每个实际整数变量证明 $Y_j\in\mathbb Z$。
\item 输出行按行序：两有限侧相等时证明 $a_j[Y]-L_j=0$；否则先证明有限下侧，
再证明有限上侧。不存在的界不产生义务，矛盾界仍须检查。
\item 最后一项永远为 $f_O[Y]-f_I\le0$，包含两个完整目标的常数。
\end{enumerate}
每项引用必须证明准确表达式与种类，缺失、多余或错误顺序均拒绝。
等式事实可以证明同表达式的不增目标；没有 Equal/None 模式。
每个叶子的一套义务始终使用一个共同 $Y$。

\section{lift 的执行协议}
\label{sec:lift-execution}
五种 lift 均先检查 witness、尺寸与资源，再在\emph{扩展前}作用域核验全部前提，
最后进入生成模型上的唯一 body。旧模型的变量、行和目标保持完整前缀，新变量目标系数为零。
新事实只在 body 可用，不能循环证明本次 lift 的前提。嵌套 lift 可用上层已经建立的坐标。
存在引理保证一个共同扩展点满足全部新条件；内核构造公式，不运行数值选择或路径搜索。

\section{完成、资源和拒绝}
完整解码先于执行，所有分支字节均须合法。普通 Absurd 只接受正的恒定 Le 或非零恒定 Eq。
每个 Block 可以 Return、Split、Absurd 或五种 lift 完成；Cases 的 Derivation 只能 Witness 或 Absurd。
没有附加性质阶段，没有可跳过的方向，也不以猜测不可达为由略过失败分支。

默认工作预算为 $10^9$、活跃事实预算为 $10^6$、组合嵌套深度为 64。
两个方向共享工作预算，局部事实分别重建。Split、Cases 和 lift body 增加深度。
超出工作、深度、位数、项数或分配限制均不接受；资源不足不是数学反例。
诊断只报告 direction、step、obligation 等已知位置，不参与决定接受。

\section{正确性不变式}
\label{sec:soundness-invariants}
对完整输入域与路径假设归纳：可见事实有效；每个分类穷尽；矛盾叶没有可行点；
每个 lift 对每个父点给出共同可行扩展；每个 Return 证明完整输出可行且目标不增。
因此每个成功方向给出一个总构造，两个方向共同推出优化等价。
这依赖精确算术、模型导入、作用域、义务生成与可信构造规则的实现正确性，
不声称 Rust 程序已由证明助手形式化验证。
